We evaluate on two topical domains: (1) primarily internal author notes and project documentation, and (2) the same domains supplemented where available by small, curated sets of published PDFs chosen to reflect mixed private/public writing scenarios and to test integration of heterogeneous evidence during section-level generation. Source materials are converted to plain text and partitioned into retrievable snippets (paragraphs, captions, section fragments) indexed with provenance metadata to enable per-section evidence injection and traceability. Because internal notes are private, we do not release raw artifacts; curated PDFs are documented in supplementary metadata where licensing permits. Evaluation uses document-holdout, generation-time retrieval, and per-section auditing (no model fine-tuning on evaluation data) to measure citation coverage and unsupported-assertion rates, and reproducibility limits and mitigations (synthetic/anonymized examples, metadata tables) are described in the supplement.